\chapter{المنحنيات المستوية}\label{ch:b1:curves}

بُني حساب \cref{ch:b1:derivative,ch:b1:taylor} من أجل
الدوال $y = f(x)$؛ وأكثر منحنيات الهندسة والميكانيك ---
المسارات، والدوائر المتدحرجة على دوائر، والمدارات --- تأبى ذلك الشكل
وتأتي بدلًا من ذلك منحنيات \emph{معلَّمية} $t \mapsto (x(t),
y(t))$ أو منحنيات \emph{قطبية} $r = r(\theta)$. ويدرس هذا الفصل
الكلاسيكي كليهما، ويُختم بالمقاطع المخروطية.

\section{المنحنيات المعلَّمية}

\begin{definition}\label{def:b1:curves:def}
\emph{المنحني المعلَّمي}\index{منحن معلَّمي} تطبيقٌ $f
\colon I \to \R^2$ مع $t \mapsto M(t) = (x(t), y(t))$، حيث $x, y$ من
الصنف $C^1$ (على الأقل) على \omterm{prop:b1:reals:intervals}{الفترة} $I$. و\emph{متجهة
السرعة} هي $f'(t) = (x'(t), y'(t))$؛ وتكون النقطة $M(t_0)$
\emph{ناظمية} إذا كان $f'(t_0) \neq (0,0)$، ويكون \emph{المماس}
هناك المستقيمَ المارّ بالمقدار $M(t_0)$ الموجَّه بالمقدار $f'(t_0)$.
\end{definition}

\begin{proof}[لماذا يوجّه $f'$ المماس]
بتايلور--يونغ مركّبةً مركّبة، $M(t) = M(t_0) + (t - t_0) f'(t_0) +
o(t - t_0)$: ومنه فاتجاه الوتر $\frac{M(t) - M(t_0)}{t - t_0}$
يؤول إلى $f'(t_0)$. وفرض الناظمية هو ما يجعل
النهاية \emph{اتجاهًا}: فإذا كان $f'(t_0) = 0$، انهارت المعروضة
إلى $M(t) = M(t_0) + o(t - t_0)$ ولم تقل شيئًا عن كيفية مغادرة
المنحني للنقطة --- ومنه المعالجة المنفصلة
للنقاط الشاذة في الطريقة أدناه، حيث تتولّى أولُ \omterm{def:b1:derivative:def}{مشتقة}
\emph{غير منعدمة} دور $f'$.
ولاحظ كذلك أن مستقيم المماس غرض هندسي: فأيّ
إعادة تعليم تغيّر $f'$ بعامل سلّمي غير معدوم و
تترك المستقيم على حاله.
\end{proof}

\begin{method}[دراسة منحن معلَّمي]\label{met:b1:curves:study}
\begin{enumerate}
  \item \emph{قلّص \omterm{def:b1:logic:sets}{مجموعة} التعريف} باستعمال التناظرات: العلاقات بين
        $M(-t)$ و $M(t + T)$ و\dots و $M(t)$ (الانعكاسات و
        الانسحابات والدورية)، ثم ارسم الجزء المقلَّص فقط.
  \item \emph{التغيرات}: جدوِل إشارتَي $x'$ و $y'$ معًا؛
        فيتحرك المنحني يمينًا/يسارًا حسب $x'$، وصعودًا/نزولًا حسب $y'$.
  \item \emph{النقاط اللافتة}: المماس الأفقي ($y' = 0 \neq
        x'$)، والمماس العمودي ($x' = 0 \neq y'$)؛ وعند
        نقطة \emph{شاذة} ($f' = 0$)، افحص المشتقات العليا
        لتجد اتجاه المماس.
  \item \emph{السلوك المقارب} عند طرفَي $I$، ثم ارسم.
\end{enumerate}
\end{method}

\begin{example}[الفروع المقاربة، معالجةً]
\label{ex:b1:curves:asymptotes}
$x(t) = t$ و $y(t) = t + \dfrac1t$ على $\intoo{0}{+\infty}$. وعندما $t
\to 0^{+}$: يكون $x \to 0$ بينما $y \to +\infty$ --- فيتسلق المنحني
على امتداد \emph{المقارب العمودي} $x = 0$ (بنهاية منتهية من أجل
إحداثية، وغير منتهية من أجل الأخرى). وعندما $t \to +\infty$:
تنفجر الإحداثيتان معًا، فاختبر مستقيمًا: يُظهر الفرق
\[
y(t) - x(t) = \frac1t \longrightarrow 0^{+}
\]
\emph{المقاربَ المائل} $y = x$، مقتربًا منه من
أعلى. وبينهما، ينعدم $y' = 1 - \frac1{t^2}$ عند $t = 1$:
فالنقطة $(1, 2)$ هي أخفض نقطة في الفرع، و $x' = 1 >
0$ في كل مكان، ومنه فالمنحني يتقدم يمينًا دائمًا. أي قوس واحد
ومقاربان وقيمة صغرى واحدة: صورةٌ كاملة من ثلاثة
حسابات --- مع ظهور الوصفة العامة تحتها:
فحين يكون $x, y \to \infty$ معًا، افحص $y/x$ بحثًا عن ميل
مرشّح (وهو هنا $\to 1$)، ثم $y - (\text{ميل})\,x$ من أجل
المقطع وجهة الاقتراب.
\end{example}

\begin{example}[النجمية]\label{ex:b1:curves:astroid}
$x(t) = \cos^3 t$ و $y(t) = \sin^3 t$.\index{النجمية} والتناظرات:
$M(t + 2\pi) = M(t)$ (فادرس دورًا واحدًا)؛ و $M(-t)$ انعكاس
$M(t)$ في محور $x$؛ و $M(\pi - t)$ في محور $y$؛ و $M(\frac\pi2 -
t)$ في القطر $y = x$: ومنه يكفي دراسة $t \in
\intcc{0}{\frac\pi4}$ ثم النشر.

والسرعة: $f'(t) = 3\sin t\cos t\,(-\cos t, \sin t)$. وعلى
$\intoo{0}{\frac\pi2}$ تكون كل النقاط ناظمية بمماس موجَّه
بالمقدار $(-\cos t, \sin t)$؛ وعند $t = 0$ (وهي النقطة $(1,0)$) تنعدم
السرعة: أي \emph{نقطة ارتداد}\index{نقطة ارتداد}، حيث ينعكس المنحني على امتداد
اتجاه المماس $(-1, 0)\cdot(\pm)$ --- وبالتناظر تقع نقاط
الارتداد الأربع عند $(\pm1, 0), (0, \pm1)$.
\end{example}

\begin{example}[منحني على شكل ثمانية، مدروسًا كاملًا]
\label{ex:b1:curves:lissajous}
$x(t) = \sin t$ و $y(t) = \sin 2t$ (وهو \emph{منحني ليساجو}).
\emph{التناظرات}: $M(t + \pi) = (-x(t), y(t))$ (انعكاس في
محور $y$)، و $M(-t) = (-x(t), -y(t))$ (تماثل مركزي)،
و $M(\pi - t) = (x(t), -y(t))$ (انعكاس في محور $x$): ومنه
يكفي دراسة $t \in \intcc{0}{\frac\pi2}$ ثم النشر.
\emph{والتغيرات}: $x' = \cos t \geq 0$ في كل مكان، بينما $y' =
2\cos 2t$ موجب قبل $t = \frac\pi4$ وسالب بعده:
فيتسلق القوس يمينًا إلى القمة
$\bigl(\frac{\sqrt2}2,\ 1\bigr)$ عند $t = \frac\pi4$ (بمماس
أفقي)، ثم ينزل يمينًا إلى $(1, 0)$ عند $t =
\frac\pi2$، حيث $x' = 0 \neq y'$: أي مماس عمودي.
\emph{والنقطة المضاعفة}: $M(0) = M(\pi) = (0,0)$، بسرعتين مختلفتين
\[
f'(0) = (1,\ 2), \qquad f'(\pi) = (-1,\ 2) :
\]
أي فرعان \emph{ناظميان} يتقاطعان عند المبدأ بزاويتين
متمايزتين --- أي نقطة مضاعفة لا نقطة شاذة: فكل مرور
ملس تمامًا، والمروران يتقاسمان موضعهما فحسب.
والمنحني كله هو الثمانية أدناه.
\end{example}

\begin{omfigure}
\begin{tikzpicture}[scale=1.5]
  \draw[->, thin] (-1.25,0) -- (1.3,0) node[below] {$x$};
  \draw[->, thin] (0,-1.25) -- (0,1.3) node[left] {$y$};
  \draw[omDef, very thick, domain=0:360, samples=200, smooth]
    plot ({sin(\x)}, {sin(2*\x)});
  \fill (0.7071, 1) circle (0.035);
  \node[above, font=\small] at (0.7071, 1)
    {$\bigl(\tfrac{\sqrt2}2, 1\bigr)$};
  \fill (0,0) circle (0.035);
\end{tikzpicture}

{\small منحني ليساجو $(\sin t, \sin 2t)$: بنقطة مضاعفة
عند المبدأ، حيث يتقاطع فرعان ناظميان بسرعتين
$(1, 2)$ و $(-1, 2)$، وبمماسات أفقية عند القمم
الأربع. وقد قلّص التناظر كل العمل إلى ربع دور.}
\end{omfigure}

\begin{omfigure}
\begin{tikzpicture}[scale=1.35]
  \draw[->, thin] (-1.3,0) -- (1.35,0) node[below] {$x$};
  \draw[->, thin] (0,-1.3) -- (0,1.35) node[left] {$y$};
  \draw[omDef, very thick, domain=0:360, samples=200, smooth]
    plot ({cos(\x)^3}, {sin(\x)^3});
  \node[below right, font=\small] at (1,0) {$1$};
  \node[above left, font=\small] at (0,1) {$1$};
\end{tikzpicture}

{\small النجمية $(\cos^3 t, \sin^3 t)$: أربعة أقواس تلتقي عند
أربع نقاط ارتداد. وهي المنحني الذي ترسمه نقطة من دائرة نصف قطرها
$\frac14$ تتدحرج داخل الدائرة الواحدية.}
\end{omfigure}

\section{المنحنيات القطبية}

\begin{definition}\label{def:b1:curves:polar}
\emph{المنحني القطبي}\index{منحن قطبي} يُعطى بالمقدار $r = r(\theta)$:
فتكون نقطة الوسيط $\theta$ هي
\[
M(\theta) = r(\theta)\,\vec u(\theta),
\qquad \vec u(\theta) = (\cos\theta, \sin\theta).
\]
ومع $\vec v(\theta) = (-\sin\theta, \cos\theta) = \vec
u\,'(\theta)$، تكون السرعة
\[
M'(\theta) = r'(\theta)\, \vec u(\theta) + r(\theta)\, \vec
v(\theta) .
\]
والنتائج: حيث $r \neq 0$، تكون النقطة ناظمية، ويصنع
المماس مع نصف القطر الزاوية $V$ المعطاة بالمقدار $\tan V =
\frac{r}{r'}$ (وهي الزاوية بين $M'$ و $\vec u$)؛ وحيث $r(\theta_0)
= 0$، يمرّ المنحني بالمبدأ \emph{بمماس هو نصف القطر
$\theta = \theta_0$} (بالاتجاه $\vec u(\theta_0)$، مقروءًا من
$M'(\theta_0) = r'(\theta_0)\vec u(\theta_0)$، أو من نهاية
الوتر: فالوتر من $O$ إلى $M(\theta)$ محمول على $\vec
u(\theta)$ نفسه، وهو يؤول إلى $\vec u(\theta_0)$ --- ومنه
تصحّ القاعدة حتى عندما $r'(\theta_0) = 0$ وتكون النقطة
شاذة، ولهذا تحصل نقاط الارتداد القطبية عند المبدأ، مثل نقطة
القلبية، على مماسها بالمجّان).
\end{definition}

\begin{example}[القلبية]\label{ex:b1:curves:cardioid}
$r(\theta) = 1 + \cos\theta$.\index{القلبية} والتناظر: $r(-\theta)
= r(\theta)$: أي انعكاس في محور $x$؛ فادرس $\theta \in
\intcc{0}{\pi}$. ويتناقص $r$ من $2$ إلى $0$؛ وعند $\theta = \pi$،
$r = 0$: فيبلغ المنحني المبدأ ماسًّا لنصف القطر
$\theta = \pi$ (وهو محور $x$)، مكوّنًا \omterm{ex:b1:curves:astroid}{نقطة ارتداد} هناك --- أي رأس
القلب. والمماس عند $\theta = 0$: $r' = 0$، ومنه $\tan V = \infty$:
أي عمودي على المحور.

وتستحق الزاوية $V$ قراءة أخرى. فعند $\theta =
\frac\pi2$: $r = 1$ و $r' = -1$، ومنه $\tan V = \frac{r}{r'} =
-1$: أي إن المماس يصنع ثلاثة أرباع الزاوية القائمة مع نصف
القطر الخارج --- فالمنحني ينثني أصلًا عائدًا نحو
نقطة ارتداده. وتعطي الصيغة $\tan V = r/r'$ اتجاهات المماس
على امتداد المنحني كله دون أيّ حساب للمقدار $M'(\theta)$
أصلًا: فهي النظير القطبي لقراءة ميل.
\end{example}

\begin{example}[من القطبي إلى الديكارتي: دائرة مستترة]
\label{ex:b1:curves:polarcircle}
ما هو \omterm{def:b1:curves:polar}{المنحني القطبي} $r = 2\cos\theta$؟ اضرب في $r$:
$r^2 = 2r\cos\theta$، أي $x^2 + y^2 = 2x$، أي
\[
(x - 1)^2 + y^2 = 1 :
\]
أي الدائرة ذات المركز $(1, 0)$ ونصف القطر $1$، المارّة بالمبدأ.
ومسك الدفاتر مهمّ: فبينما يجري $\theta$ على
$\intoc{-\frac\pi2}{\frac\pi2}$ تُمسح الدائرة كلها
مرة واحدة بالضبط (لأن $r$ ينعدم عند الطرفين)، وعند $\theta =
\pm\frac\pi2$ تعطي القاعدة «المماس عند المبدأ على امتداد نصف القطر
$\theta = \theta_0$» مماسًا \emph{عموديًا} هناك ---
مطابقًا للهندسة، لأن المحور العمودي ماسٌّ فعلًا
لهذه الدائرة عند $O$. وأمّا من أجل $\theta$ بعد ذلك المجال، فإن $r < 0$
يعيد رسم الدائرة نفسها: وهو تذكير بأن \omterm{def:b1:curves:polar}{المنحني القطبي} غرضٌ
\emph{معلَّمي}، مسموحٌ له بالمرور فوق نفسه.
\end{example}

\begin{omfigure}
\begin{tikzpicture}[scale=1.5]
  \draw[->, thin] (-1.2,0) -- (1.3,0) node[below] {$x$};
  \draw[->, thin] (0,-1.2) -- (0,1.3) node[left] {$y$};
  \draw[omProp, very thick, domain=0:360, samples=300, smooth]
    plot ({cos(2*\x)*cos(\x)}, {cos(2*\x)*sin(\x)});
  \draw[thin, dashed] (-0.85,-0.85) -- (0.85,0.85);
  \draw[thin, dashed] (-0.85,0.85) -- (0.85,-0.85);
\end{tikzpicture}

{\small الوردة ذات البتلات الأربع $r = \cos 2\theta$
(\cref{exo:b1:curves:4}). وتُرسم البتلتان على امتداد محور $y$
بالمقدار $r < 0$ (فتُرسم النقطة على نصف القطر المقابل
للمقدار $\theta$)؛ والقطران المتقطّعان $\theta = \pm\frac\pi4$ هما
المماسان عند المبدأ، حيث ينعدم $r$.}
\end{omfigure}

\begin{omfigure}
\begin{tikzpicture}[scale=1.15]
  \draw[->, thin] (-0.7,0) -- (2.4,0) node[below] {$x$};
  \draw[->, thin] (0,-1.7) -- (0,1.7) node[left] {$y$};
  \draw[omThm, very thick, domain=0:360, samples=200, smooth]
    plot ({(1+cos(\x))*cos(\x)}, {(1+cos(\x))*sin(\x)});
  \node[below, font=\small] at (2,0) {$2$};
\end{tikzpicture}

{\small القلبية $r = 1 + \cos\theta$. وتُقرأ المنحنيات القطبية
بمسح الزاوية: فينتفخ نصف القطر ويتقلص كلما دار $\theta$.}
\end{omfigure}

\section{طول القوس}

\begin{definition}[طول قوس]\label{def:b1:curves:arclength}
\emph{طول}\index{طول القوس} قوس من الصنف $C^1$ $f \colon
\intcc{a}{b} \to \R^2$ هو تكامل السرعة:
\[
L = \int_a^b \norm{f'(t)}\,\dd t
= \int_a^b \sqrt{x'(t)^2 + y'(t)^2}\;\dd t .
\]
(والتحفيز: على \omterm{prop:b1:reals:intervals}{فترة} صغيرة، $M(t + h) \approx M(t) + h
f'(t)$، ومنه فالقوس قريب من مضلع تُجمع أطوال قطعه
لتعطي مجموع ريمان للمقدار $\norm{f'}$،
\cref{thm:b1:integration:riemann}.) ومن أجل \omterm{def:b1:curves:polar}{منحن قطبي} $r =
r(\theta)$، يكون للسرعة $r'\vec u + r\vec v$ مركّبتان
متعامدتان، ومنه
\[
L = \int_{\theta_1}^{\theta_2}
\sqrt{r'(\theta)^2 + r(\theta)^2}\;\dd\theta .
\]
ولا يتعلق الطول بالتعليم المختار (الرتيب، من الصنف
$C^1$): إذ إن التعويض $t = \varphi(s)$ في
\omterm{thm:b1:integration:def}{التكامل} (\cref{thm:b1:integration:parts}) يضرب $f'$ في
$\varphi'$ و $\dd t$ في $\varphi'^{-1}$.
\end{definition}

\begin{example}[تحقق من السلامة: الدائرة]\label{ex:b1:curves:circlelength}
من أجل $f(t) = (R\cos t, R\sin t)$ على $\intcc{0}{2\pi}$: $\norm{f'} =
R$، ومنه $L = 2\pi R$ --- فيعيد التعريف المحيط.
واختبار إعادة التعليم: يرسم \omterm{def:b1:logic:map}{التطبيق} $g(t) = (R\cos 2t, R\sin 2t)$ على
$\intcc{0}{\pi}$ الدائرةَ \emph{نفسها} بسرعة مضاعفة
$\norm{g'} = 2R$، و
\[
\int_0^{\pi} 2R\,\dd t = 2\pi R
\]
مرة أخرى: نصف الزمن وضِعف السرعة والطول نفسه --- وهو
الحفظ الموعود به في التعريف، مراقَبًا مرة على أعداد.
(وتشغيل $g$ على $\intcc{0}{2\pi}$ كله سيعطي $4\pi R$:
فالمنحني المقطوع مرتين أطول مرتين \emph{بوصفه رحلة}؛
فالطول يقيس المسار المعلَّم، ومسك الدفاتر الأمين
\omterm{prop:b1:reals:intervals}{للفترة} جزء من الحساب.)
ومن أجل النجمية $(\cos^3t, \sin^3t)$: $\norm{f'} = 3\abs{\sin
t\cos t} = \tfrac32\abs{\sin 2t}$، وبالتناظر $L = 4\int_0^{\pi/2}
\tfrac32\sin 2t\,\dd t = 6$: أي منحني مرسوم داخل الدائرة الواحدية،
طوله $6 < 2\pi$. وتقيس مسألة نهاية الأسبوع أشهر
قوس على الإطلاق.
\end{example}

\begin{omfigure}
\begin{tikzpicture}[scale=0.85]
  \draw[->, thin] (-0.5,0) -- (7.2,0) node[below] {$x$};
  \draw[->, thin] (0,-0.4) -- (0,2.6) node[left] {$y$};
  \draw[omThm, very thick, domain=0:6.2832, samples=160, smooth]
    plot ({\x - sin(\x r)}, {1 - cos(\x r)});
  \draw[omDef, thick] (2,1) circle (1);
  \fill[omDef] (2,1) circle (0.04);
  \fill (2 - 0.9093, 1.4161) circle (0.06);
  \draw[omProp, thick] (2,0) -- (2 - 0.9093, 1.4161) -- (2,2);
  \fill (2,0) circle (0.05);
  \fill (2,2) circle (0.05);
  \node[below, font=\small] at (2,0) {$C$};
  \node[above, font=\small] at (2,2) {$T$};
  \node[left, font=\small] at (2 - 0.9093, 1.4161) {$M$};
  \node[below, font=\small] at (6.2832,0) {$2\pi R$};
\end{tikzpicture}

{\small قوس واحد من الدويرية (بنصف قطر $R = 1$)، والدائرة
المتدحرجة عند $t = 2$، ومستقيما الإرشاد في مسألة نهاية الأسبوع:
فالوتر $MC$ إلى نقطة التماس \emph{ناظم} على
المنحني، والوتر $MT$ إلى قمة الدائرة \emph{مماس}.}
\end{omfigure}

\section{المقاطع المخروطية}

\begin{definition}[التعريف بالبؤرة والدليل]\label{def:b1:curves:conics}
ثبّت نقطة $F$ (وهي \emph{البؤرة})، ومستقيمًا $D$ لا يمرّ بالمقدار $F$
(وهو \emph{الدليل})، و $e > 0$ (وهو \emph{الانحراف المركزي}). ويكون
\emph{المقطع المخروطي}\index{مقطع مخروطي} لهذه المعطيات هو
\[
\mathcal{C} = \{M : d(M, F) = e\; d(M, D)\}:
\]
أي \emph{قطع ناقص} من أجل $e < 1$، و\emph{قطع مكافئ} من أجل $e = 1$، و
\emph{قطع زائد} من أجل $e > 1$. (وتظهر الدائرة نهايةً منحلّة
$e \to 0$.)
\end{definition}

\begin{example}[التعريف، متحققًا منه على قطع مكافئ]
\label{ex:b1:curves:parabolacheck}
خذ القطع المكافئ $y^2 = 4x$، أي $2p = 4$: فالبؤرة $F = (1, 0)$
والدليل $D : x = -1$ (وتضعهما الصورة المختزلة
\cref{thm:b1:curves:reduced} عند $\pm\frac p2$). وعند
النقطة $M = (1, 2)$ من المنحني:
\[
MF = \sqrt{(1-1)^2 + 2^2} = 2,
\qquad
d(M, D) = 1 - (-1) = 2 :
\]
فهما متساويان، كما يقتضي $e = 1$. وعند $M' = (4, 4)$: $MF' = \sqrt{9 +
16} = 5$ و $d(M', D) = 5$ مرة أخرى. فالتعريف بالبؤرة والدليل
ليس تجريدًا: بل هو زوج مسافات يمكن
قياسه على أيّ نقطة، وجبر المعادلات
المختزلة ليس إلا هذا القياس مصنوعًا مرة واحدة و
إلى الأبد.
\end{example}

\begin{theorem}[المعادلات المختزلة]\label{thm:b1:curves:reduced}
في معلم متعامد ممنظم محسن الاختيار:
\[
\text{قطع ناقص: } \frac{x^2}{a^2} + \frac{y^2}{b^2} = 1
\quad (a \geq b > 0),
\qquad
\text{قطع زائد: } \frac{x^2}{a^2} - \frac{y^2}{b^2} = 1,
\qquad
\text{قطع مكافئ: } y^2 = 2px,
\]
مع، من أجل القطع الناقص: بؤرتين عند $(\pm c, 0)$ و $c = \sqrt{a^2 - b^2}$ و
$e = \frac ca$، والتمييز ثنائي البؤرة $MF + MF' = 2a$؛
ومن أجل القطع الزائد: $c = \sqrt{a^2 + b^2}$ و $e = \frac ca$ و
$\abs{MF - MF'} = 2a$ ومقاربان $y = \pm\frac ba x$.
\end{theorem}

\begin{proof}
خذ البؤرة عند المبدأ والدليل عموديًا، $x = -h$
($h > 0$): فيُقرأ الشرط $MF^2 = e^2 d(M, D)^2$ هكذا $x^2 + y^2 =
e^2 (x + h)^2$. ومن أجل $e = 1$: $y^2 = 2hx + h^2$، أي قطع مكافئ بعد
الانسحاب $x \mapsto x - \frac h2$ (ومنه $p = h$). ومن أجل $e \neq 1$:
بإتمام المربع في $x$،
\[
(1 - e^2)\Bigl(x - \frac{e^2 h}{1 - e^2}\Bigr)^{\!2} + y^2
= e^2h^2 + \frac{e^4 h^2}{1 - e^2}
= \frac{e^2 h^2}{1 - e^2} .
\]
ضع $X = x - \frac{e^2h}{1-e^2}$ (وهو انسحاب للمعلم). وعندما $e < 1$،
اقسم على الطرف الأيمن الموجب:
$\frac{X^2}{a^2} + \frac{y^2}{b^2} = 1$ مع $a = \frac{eh}{1 -
e^2}$ و $b = \frac{eh}{\sqrt{1-e^2}}$؛ وعندما $e > 1$، ينقلب طرفا
المعروضة انقلابًا مناسبًا وتعطي القسمة نفسها
$\frac{X^2}{a^2} - \frac{y^2}{b^2} = 1$ مع $a = \frac{eh}{e^2 -
1}$ و $b = \frac{eh}{\sqrt{e^2 - 1}}$. وتتبع قيم $c$ المذكورة
(لأن $c^2 = a^2 - b^2$ أو $a^2 + b^2$ يعطي $c = ea$ في
الحالتين، فيضع البؤرة في موضعها الصحيح)، وأمّا الخصائص ثنائية البؤرة فهي
تحقّقات مباشرة على المعادلات المختزلة. وبالتفصيل من أجل
القطع الناقص، مع $F' = (c, 0)$ و $M = (X, y)$ على المنحني:
\[
MF'^2 = (X - c)^2 + y^2
= X^2 - 2cX + c^2 + b^2\Bigl(1 - \frac{X^2}{a^2}\Bigr)
= \frac{c^2}{a^2}X^2 - 2cX + a^2
= (a - eX)^2 ,
\]
باستعمال $c^2 + b^2 = a^2$ و $e = \frac ca$؛ ولأن $\abs X \leq
a$ و $e < 1$، يكون $a - eX > 0$، ومنه $MF' = a - eX$ دون استخراج أيّ جذر
تربيعي قط. ويعطي الحساب المرآتي $MF = a + eX$،
ومنه $MF + MF' = 2a$، وهو ثابت. وأمّا من أجل القطع الزائد فيعطي الجبر
نفسه $MF = \abs{a + eX}$ و $MF' = \abs{eX - a}$،
وفرقهما $\pm2a$ تبعًا للفرع.
\end{proof}

\begin{example}\label{ex:b1:curves:conicexample}
$\frac{x^2}{25} + \frac{y^2}{9} = 1$: قطع ناقص، $a = 5$ و $b = 3$ و $c
= 4$: فالبؤرتان $(\pm4, 0)$، والانحراف المركزي $\frac45$. وتعليمه:
$(5\cos t, 3\sin t)$ --- أي دائرة ممطوطة تمطيطًا لا تناظريًا؛ و
مجموع المسافتين إلى البؤرتين من أيّ نقطة له هو $10$.
\end{example}

\begin{example}[قراءة مدار من معادلته القطبية]
\label{ex:b1:curves:orbit}
\omterm{def:b1:curves:conics}{المقطع المخروطي} القطبي $r = \dfrac{1}{1 + \frac12\cos\theta}$
(\cref{exo:b1:curves:7} مع $p = 1$ و $e = \frac12$) قطعٌ
ناقص \emph{ببؤرة عند المبدأ} --- أي هندسة
مدار كوكبي والشمس عند $O$. واستخرج كل شيء من
$p$ و $e$: فمن اختزال
\cref{thm:b1:curves:reduced}، $a = \dfrac{p}{1 - e^2} =
\dfrac{1}{3/4} = \dfrac43$ و $c = ea = \dfrac23$. ويتحقق الحضيضان
منه دون أيّ نظرية أصلًا:
\[
r(0) = \frac{1}{3/2} = \frac23 = a - c
\quad (\text{الحضيض}),
\qquad
r(\pi) = \frac{1}{1/2} = 2 = a + c
\quad (\text{الأوج}),
\]
ومجموعهما $\frac23 + 2 = \frac83 = 2a$ يستعيد المحور
الأكبر. والصورة القطبية هي الطبيعية كلما تميّزت بؤرة
فيزيائيًا؛ والصورة الديكارتية المختزلة كلما تميّزت
محاور التناظر. والتحويل بينهما هو بالضبط
ما يفعله حساب إتمام المربع في المبرهنة.

والانحراف المركزي مؤشرًا، أخيرًا: أبقِ $p = 1$ وأدر $e$ في $r
= \frac{1}{1 + e\cos\theta}$. فعند $e = 0$: الدائرة $r = 1$. وعند
$e = \frac12$: القطع الناقص المدروس للتوّ، و $r$ يتذبذب بين
$\frac23$ و $2$. وعند $e = 1$: $r(\theta) \to \infty$ عندما $\theta
\to \pi$ --- فلا ينغلق المنحني بعد: أي قطع مكافئ، أُبعدت
نقطته الأبعد إلى اللانهاية. وعند $e = 2$: ينعدم المقام
عند $\cos\theta = -\frac12$، ولا ينجو إلا $\theta \in
\intoo{-\frac{2\pi}3}{\frac{2\pi}3}$: أي فرع واحد من قطع
زائد، يفلت على امتداد اتجاهين مقاربين. أي صيغة
واحدة، وعائلة المقاطع المخروطية كلها، ونقطتا الانتقال $e =
1$ مرئيتان لحظةَ يبلغ المقام الصفرَ أول مرة.
\end{example}

\begin{remark}[مزالق شائعة]
\label{rem:b1:curves:pitfalls}
\emph{النقاط المضاعفة ليست نقاطًا شاذة}: فعند
تقاطع ذاتي، يكون كل فرع ناظميًا؛ و«الشاذة» تشير إلى
$f'(t_0) = 0$ من أجل قيمة واحدة للوسيط
(\cref{ex:b1:curves:lissajous} إزاء نقاط ارتداد النجمية).
\emph{المماس العمودي إزاء \omterm{ex:b1:curves:astroid}{نقطة الارتداد}}: فالمقدار $x'(t_0) = 0 \neq y'(t_0)$
نقطة ناظمية بمماس عمودي؛ ولا يقتضي النشرَ من الرتبة العليا
إلا $x' = y' = 0$.
\emph{أنصاف الأقطار السالبة تُرسم على نصف القطر المقابل}: فمن أجل $r(\theta) <
0$ تكون النقطة $-\abs{r}\,\vec u(\theta)$، بالزاوية $\theta +
\pi$ (\cref{exo:b1:curves:5})؛ ونسيان هذا يضيّع الحلقات الداخلية
أو يضاعف المنحنيات.
\emph{\omterm{def:b1:curves:arclength}{طول القوس} يكامل $\norm{f'}$ لا $f'$}: فاقسم
\omterm{thm:b1:integration:def}{التكامل} عند جذور السرعة --- فمن أجل النجمية،
تكامل $\frac32\sin 2t$ على دور كامل دون قيم مطلقة
يعطي $0$ لا $6$.
\emph{يجب اختزال المقاطع المخروطية قبل قراءتها}: ففي $9x^2 +
25y^2 - 36x - 50y - 164 = 0$، لا تُرى المحاور ولا المركز
ولا الانحراف المركزي حتى تُتمّ المربعات
(\cref{exo:b1:curves:6})؛ وانعدام الجزء من الدرجة الثانية على
متغير واحد يعني قطعًا مكافئًا، لا «قطعًا ناقصًا منحلًّا».
\end{remark}

\begin{remark}[إلى أين تذهب هذه المنحنيات]
\label{rem:b1:curves:whereused}
المنحنيات المعلَّمية لغة الميكانيك: فالمسارات
منحنيات، ومتجهات السرعة متجهات مماسة، ويكون طول
قوس \cref{def:b1:curves:arclength} المسافةَ المقطوعة. و
تعود المقاطع المخروطية حيثما فعل قانون التربيع العكسي --- فالمدارات
الكوكبية قطوع ناقصة والشمس في بؤرة. وفي
\cref{ch:b1:multivar}، تصير المنحنيات مجموعات المستوى لدوال
في متغيرين، ويلتقي مماس هذا الفصل
بتدرج الفصل التالي. ويضيف مجلّد السنة~الثانية الانحناء و
الصورة القانونية الموضعية لمنحن؛ وتستخرج مسألة نهاية الأسبوع أدناه
أصلًا، بأدوات هذه السنة وحدها، كلَّ ما عرفه
القرن السابع عشر عن أشهر منحنياته.
\end{remark}

\begin{remark}[منظورات داخل الكتاب 3]
\label{rem:b1:curves:perspectives}
يستهلك هذا الفصل النصف الأول كله من المجلّد و
يغذّي الفصل الأخير. \emph{المستهلَك}: فمتجهات المماس
مشتقات (\cref{ch:b1:derivative})، \omterm{def:b1:curves:arclength}{وطول القوس} والمساحات
تكاملات (\cref{ch:b1:integration})، وتُحسم نقاط الارتداد
بنشور تايلور (\cref{ch:b1:taylor})، \omterm{pb:b1:curves:1}{ومتساوية الزمن}
معادلة تفاضلية خطية (\cref{ch:b1:diffeq})، وكل
مسافة وزاوية إقليدية (\cref{ch:b1:euclid}).
\emph{والمغذَّى}: ففي \cref{ch:b1:multivar}، يُشتق منحن $t \mapsto (x(t),
y(t))$ مرسوم داخل \omterm{def:b1:logic:sets}{مجموعة} مستوى $\{f = c\}$ بقاعدة
السلسلة، وتزوّج المتطابقة الناتجة $\langle \nabla f,\
f'(t)\rangle = 0$ متجهاتِ المماس في هذا الفصل بتدرجات
الفصل التالي --- فمستقيم مماس المنحني و
الاتجاه الناظم للسطح هما الحساب نفسه منظورًا إليه
من الضفتين.
\end{remark}

\section{تمارين}

\begin{exercise}[$\star$]\label{exo:b1:curves:1}
ادرس وارسم المنحني $x(t) = t^2$ و $y(t) = t^3$ (التناظرات و
التغيرات والسلوك عند النقطة الشاذة $t = 0$).
\end{exercise}

\begin{exercise}[$\star$]\label{exo:b1:curves:2}
من أجل الدويرية $x(t) = t - \sin t$ و $y(t) = 1 - \cos t$ (وهي
مسار نقطة من عجلة متدحرجة نصف قطرها $1$): عيّن
تناظر الانسحاب بالدور، والنقاط الشاذة، و
اتجاه المماس عند $t = 0$ \emph{(انشر $x$ و $y$ حتى أوائل
الرتب غير المعدومة)}.
\end{exercise}

\begin{exercise}[$\star$]\label{exo:b1:curves:3}
أعطِ مستقيم المماس للمنحني $(\cos^3 t, \sin^3 t)$ عند $t =
\frac\pi4$، وتحقق من أن قطعة هذا المماس التي يقطعها
المحوران طولها $1$ --- وهي خاصية شهيرة للنجمية (صحيحة عند
كل نقطة ناظمية).
\end{exercise}

\begin{exercise}[$\star$]\label{exo:b1:curves:4}
ارسم المنحنيين القطبيين $r = \cos 2\theta$ (الوردة ذات البتلات الأربع) و
$r = \frac{1}{\cos\theta}$ على $\intoo{-\frac\pi2}{\frac\pi2}$
(وتعرّف على مستقيم).
\end{exercise}

\begin{exercise}[$\star\star$]\label{exo:b1:curves:5}
ادرس \omterm{def:b1:curves:polar}{المنحني القطبي} $r = 1 + 2\cos\theta$ (وهو \emph{حلزونة}): \omterm{def:b1:logic:sets}{مجموعة} التعريف
حيث $r \geq 0$ إزاء $r < 0$ (فالنقاط المرسومة بنصف قطر سالب
تقع على نصف القطر المقابل)، والمرور بالمبدأ، \omterm{def:b1:structures:ring}{والحلقة} الداخلية،
والرسم.
\end{exercise}

\begin{exercise}[$\star\star$]\label{exo:b1:curves:6}
عيّن \omterm{def:b1:curves:conics}{المقطع المخروطي} $9x^2 + 25y^2 - 36x - 50y - 164 = 0$: اختزل بإتمام
المربعات، وأعطِ المركز وأنصاف المحاور والبؤرتين والانحراف المركزي.
\end{exercise}

\begin{exercise}[$\star\star$]\label{exo:b1:curves:7}
برهن على أن المعادلة القطبية \omterm{def:b1:curves:conics}{لمقطع مخروطي} ببؤرة عند المبدأ
هي
\[
r = \frac{p}{1 + e\cos\theta}
\]
(بالدليل عموديًا على بعد $\frac pe$ يمين البؤرة).
وأيّ قيم $\theta$ مسموح بها عندما $e > 1$؟
\end{exercise}

\begin{exercise}[$\star\star$]\label{exo:b1:curves:8}
من الخاصية ثنائية البؤرة $MF + MF' = 2a$، استنتج «بناء
البستاني» للقطع الناقص، وبرهن على أن المماس عند $M$
يصنع زاويتين متساويتين مع $MF$ و $MF'$ \emph{(خاصية الانعكاس؛
اشتق $\norm{M(t) - F} + \norm{M(t) - F'} = 2a$ و
فسّر انعدام مجموع جداءات المتجهات الواحدية)}.
\end{exercise}

\begin{exercise}[$\star\star\star$]\label{exo:b1:curves:9}
\emph{لمنيسكة برنولي} هي \omterm{def:b1:curves:polar}{المنحني القطبي} $r^2 = \cos
2\theta$ (خذ $r = \sqrt{\cos2\theta}$ حيث تكون معرَّفة).
\begin{enumerate}
  \item أعطِ \omterm{def:b1:logic:sets}{مجموعة} تعريفها وتناظراتها ومماساتها عند المبدأ، و
        ارسمها.
  \item برهن على أنها المحلّ الهندسي للنقاط $M$ التي تحقق $MF \cdot MF'
        = \frac12$ حيث $F, F' = \bigl(\pm\frac{1}{\sqrt2},
        0\bigr)$. \emph{(احسب $MF^2\,MF'^2$ \omterm{prop:b1:vspaces:coordinates}{بالإحداثيات}
        القطبية.)}
\end{enumerate}
\end{exercise}

\begin{exercise}[$\star\star$]\label{exo:b1:curves:10}
احسب الطول الكلي للقلبية $r = 1 + \cos\theta$
باستعمال الصيغة القطبية في \cref{def:b1:curves:arclength}
\emph{(فالمتطابقة $1 + \cos\theta = 2\cos^2\frac\theta2$ تحوّل
الجذر التربيعي إلى $2\abs{\cos\frac\theta2}$)}.
\end{exercise}

\begin{exercise}[$\star\star$]\label{exo:b1:curves:11}
بيّن أن مماس القطع الناقص $(a\cos t, b\sin t)$ عند
نقطة الوسيط $t$ معادلته
\[
\frac{x\cos t}{a} + \frac{y\sin t}{b} = 1 ,
\]
واستنتج مماس $\frac{x^2}{a^2} + \frac{y^2}{b^2} = 1$
عند نقطة $(x_0, y_0)$ من القطع الناقص: $\frac{x\,x_0}{a^2} +
\frac{y\,y_0}{b^2} = 1$ (وهي قاعدة «فصل المربعات»).
\end{exercise}

\begin{exercise}[$\star\star\star$]\label{exo:b1:curves:12}
\emph{الحلزون اللوغاريتمي} هو \omterm{def:b1:curves:polar}{المنحني القطبي} $r =
\eu^{k\theta}$ ($k > 0$ مثبَّت و $\theta \in \R$).
\begin{enumerate}
  \item بيّن أن الزاوية $V$ بين نصف القطر و
        المماس \emph{ثابتة} ($\tan V = \frac1k$) --- فيقطع
        الحلزون كل نصف قطر بالزاوية نفسها.
  \item بيّن أن تدوير الحلزون بزاوية $c$ يرسله
        على تسليمه بالعامل $\eu^{-kc}$: فكل
        دوران للحلزون تكبيرٌ له
        (تشابه ذاتي).
  \item احسب \omterm{def:b1:curves:arclength}{طول القوس} $\theta \in
        \intoc{-\infty}{\theta_0}$ (نهايةً لأطوال على
        $\intcc{A}{\theta_0}$ عندما $A \to -\infty$) ولاحظ
        أنه منته: أي منحني يلتفّ حول المبدأ عددًا لا يُحصى
        من المرات، وطوله منته.
\end{enumerate}
\end{exercise}

\section{مسألة: الدويرية، ملكة المنحنيات}

\begin{problem}\label{pb:b1:curves:1}
عجلة نصف قطرها $R$ تتدحرج دون انزلاق على امتداد محور $x$؛
والنقطة من الإطار التي كانت في البداية عند المبدأ ترسم
\emph{الدويرية}. وقد تصارع القرن السابع عشر على هذا المنحني
--- فوزن غاليليو قصاصات ورقية له، وقاسه رن،
وحسب روبرفال مساحته، وبنى هويغنس ساعات عليه ---
وكل نتيجة من نتائجهم في متناول هذا الفصل. ونبرهن
على الكلاسيكيات الأربع: بناء المماس، وطول رن
$8R$، والمساحة $3\pi R^2$، وخاصية هويغنس في
تساوي الزمن.\index{الدويرية}\index{متساوية الزمن}

\medskip
\noindent\textbf{الجزء 1 --- التدحرج والمماس.}
\begin{enumerate}
  \item بعد أن تدور العجلة بزاوية $t$، يقع مركزها
        عند $\Omega(t) = (Rt, R)$ (بالتدحرج دون انزلاق:
        فمسافة التماس $=$ القوس المتدحرج). بيّن أن النقطة
        المعلَّمة تقع عند
        \[
        M(t) = \bigl(R(t - \sin t),\; R(1 - \cos t)\bigr).
        \]
  \item احسب $f'(t)$ وبيّن $\norm{f'(t)} =
        2R\,\abs{\sin\frac t2}$؛ وحدّد النقاط الشاذة
        (نقاط الارتداد --- قارن \cref{exo:b1:curves:2}) وقمة
        كل قوس.
  \item ليكن $C(t) = (Rt, 0)$ نقطة التماس ولتكن $T(t) =
        (Rt, 2R)$ قمة العجلة. برهن على أنه من أجل $0 < t <
        2\pi$ تكون المتجهة $M - C$ \emph{ناظمة} على المنحني عند
        $M(t)$ وتكون المتجهة $T - M$ \emph{مماسة}: أي إن رسم
        المماس لدويرية يكون بوصل النقطة بقمة
        دائرتها المتدحرجة. \emph{(عمّل كل شيء عبر
        $\sin\frac t2$ و $\cos\frac t2$.)}
  \item فسّر السؤال 3 حركيًا: فنقطة التماس هي
        \emph{المركز اللحظي للدوران}، وتساوي سرعة $M$
        مسافتها إلى $C$ (من أجل سرعة زاوية واحدية).
        وتحقق من $\norm{M - C} = 2R\abs{\sin\frac t2}$.
  \item برهن على علاقة الارتفاع والسرعة
        \[
        \norm{f'(t)}^2 = 2R\,y(t) :
        \]
        على دويرية مقطوعة بسرعة زاوية واحدية، تكون
        السرعة عند كل نقطة بالضبط سرعةَ السقوط الحر من أجل
        هبوط يساوي الارتفاع الحالي. (احتفظ بهذا من أجل
        الجزء 4.)
\end{enumerate}

\medskip
\noindent\textbf{الجزء 2 --- مبرهنة رن: \omterm{def:b1:curves:arclength}{طول القوس}
$8R$.}
\begin{enumerate}[resume]
  \item باستعمال \cref{def:b1:curves:arclength}، احسب طول
        قوس واحد:
        \[
        L = \int_0^{2\pi} 2R\sin\frac t2\,\dd t = 8R
        \]
        (وهي مبرهنة رن، 1658). أي أربعة أقطار عجلة، ودون أيّ
        $\pi$ في أيّ موضع.
  \item احسب \omterm{def:b1:curves:arclength}{طول القوس} من \omterm{ex:b1:curves:astroid}{نقطة الارتداد}: $s(t) =
        4R\bigl(1 - \cos\frac t2\bigr)$، وتحقق من $s(2\pi) =
        8R$.
  \item والآن قِس القوس من \emph{القمة} $t = \pi$:
        $\sigma(t) = \abs{4R\cos\frac t2}$. برهن على العلاقة
        الذاتية
        \[
        \sigma^2 = 8R\,\bigl(2R - y\bigr) :
        \]
        أي إن مربع المسافة القوسية من القمة يتناسب مع
        الهبوط في الارتفاع تحت القمة.
  \item وتحققات السلامة: استعد من السؤال 8 أن نصف القوس
        من القمة إلى \omterm{ex:b1:curves:astroid}{نقطة الارتداد} طوله $4R$، وقارن بحساب
        النجمية في \cref{ex:b1:curves:circlelength}
        --- فلكلا المنحنيين أطوال جبرية خالية من $\pi$؛
        وفسّر ما يجعل هذا ممكنًا مع أن كليهما
        مبنيّ من دوائر. \emph{(انظر في شكل
        $\norm{f'}$.)}
\end{enumerate}

\medskip
\noindent\textbf{الجزء 3 --- مساحة روبرفال: $3\pi R^2$.}
\begin{enumerate}[resume]
  \item سوّغ أن المساحة بين قوس واحد والأرض هي
        $A = \int_0^{2\pi} y(t)\,x'(t)\,\dd t$ \emph{(بالتعويض
        $x = x(t)$ في $\int y\,\dd x$،
        \cref{thm:b1:integration:parts}؛ و $x$ متزايد)}.
  \item احسب
        \[
        A = R^2\int_0^{2\pi}(1 - \cos t)^2\,\dd t = 3\pi R^2 :
        \]
        أي \emph{ثلاثة} أضعاف مساحة العجلة بالضبط ---
        وهي النسبة التي خمّنها غاليليو بالوزن.
  \item والطريقة نفسها من أجل النجمية $(\cos^3t, \sin^3t)$: بيّن
        أن المساحة المحاطة هي $\frac{3\pi}8$
        \emph{(خطّط $\sin^4 t\cos^2 t$؛ فلا ينجو إلا الحدّ
        الثابت على دور كامل)}.
  \item تحقق من سلامة صيغة السؤال 10 على نصف الدائرة الواحدية
        العلوي $(\cos t, \sin t)$، مع $t$ من $\pi$ إلى $0$:
        فهل تعيد $\frac\pi2$؟
\end{enumerate}

\medskip
\noindent\textbf{الجزء 4 --- متساوية هويغنس الزمنية.} اقلب
القوس: تنزلق خرزة بلا احتكاك، تحت الثقالة $g$، داخل
الوعاء الدويري
\[
x(t) = R(t + \sin t), \qquad y(t) = R(1 - \cos t)
\qquad (t \in \intcc{-\pi}{\pi}),
\]
وأخفض نقطة فيه هي المبدأ ($y$ مقيسًا صعودًا).
\begin{enumerate}[resume]
  \item احسب السرعة $\norm{f'(t)} = 2R\cos\frac t2$، \omterm{def:b1:curves:arclength}{وطول
        القوس} من الأسفل $s(t) = 4R\sin\frac t2$، و
        برهن على المتطابقة المفتاحية
        \[
        y = \frac{s^2}{8R} .
        \]
  \item تخضع الخرزة المطلقة من السكون من نقطة الوسيط
        $t_0 > 0$ لحفظ الطاقة: فإذا رمز $s(\tau)$
        إلى موضعها القوسي عند الزمن $\tau$، فإن $\frac12\bigl(
        \frac{\dd s}{\dd\tau}\bigr)^2 + g\,y = g\,y_0$. أعد كتابة
        هذا، باستعمال السؤال 14، في صورة
        \[
        \Bigl(\frac{\dd s}{\dd\tau}\Bigr)^{\!2}
        = \frac{g}{4R}\,\bigl(s_0^2 - s^2\bigr),
        \qquad s_0 = s(t_0).
        \]
  \item اشتق بالنسبة إلى $\tau$ واحصل على المهتزّ
        التوافقي
        \[
        \frac{\dd^2 s}{\dd\tau^2} = -\frac{g}{4R}\,s ;
        \]
        وحُلَّه بالمقدار \cref{thm:b1:diffeq:homogeneous2} وبالشروط
        الابتدائية: $s(\tau) = s_0\cos(\omega\tau)$ و
        $\omega = \sqrt{g/4R}$.
  \item استنتج \emph{خاصية تساوي الزمن} (هويغنس، 1659):
        فزمن بلوغ الأسفل،
        \[
        T_{\downarrow} = \frac{\pi}{2\omega}
        = \pi\sqrt{\frac Rg}\,,
        \]
        لا يتعلق بنقطة الإطلاق --- فالخرزات المطلقة
        معًا من أيّ ارتفاعين في الوعاء تصل
        معًا.
  \item تحقق بالتعويض من أن $s(\tau) =
        s_0\cos\omega\tau$ يحقق معادلة الطاقة من الرتبة الأولى
        في السؤال 15 تحقيقًا مضبوطًا (لا \omterm{def:b1:derivative:def}{المشتقة}
        منها فقط)، وفسّر في جملة واحدة لماذا يكون البندول
        الدائري متساوي الزمن \emph{تقريبًا} فقط بينما
        تكون الدويرية كذلك بالضبط.
  \item وعدديًا: أيّ نصف قطر $R$ يجعل زمن النزول
        ثانيةً واحدة بالضبط ($g = 9.81$)؟ ولاحظ كم يقترب
        الجواب من المتر الواحد، وكيف يقارن دور
        الاهتزاز الكامل $2\pi\sqrt{4R/g}$ بصيغة بندول
        الزوايا الصغيرة $2\pi\sqrt{\ell/g}$ من أجل $\ell = 4R$.
\end{enumerate}

\medskip
\noindent\textbf{الجزء 5 --- الأرباح، والتوليفة.}
\begin{enumerate}[resume]
  \item طبّق قاعدة المماس في السؤال 3 عند $t =
        \frac\pi2$ (بأخذ $R = 1$): احسب $M$ و $T$ و
        اتجاه $MT$، وتحقق منه إزاء
        $f'(\frac\pi2)$.
  \item (التروكويدات) علّم بدل ذلك نقطة على بعد $d$ من
        المركز ($d \neq R$): فيكون المنحني $x = Rt -
        d\sin t$ و $y = R - d\cos t$. بيّن أنه من أجل $d < R$ يكون
        المنحني ناظميًا في كل مكان وبلا أيّ سلوك شاذ
        ($x' > 0$: فهو يتقدم)، بينما من أجل $d
        > R$ تغيّر الفاصلة $x'$ إشارتها ويصنع المنحني
        حلقات --- أي عجلة السكة ذات \omterm{def:b1:topology:closure}{الحافة} التي تسافر نقاط
        إطارها \emph{إلى الوراء}.
  \item (تشويق أقصر زمن) من \omterm{ex:b1:curves:astroid}{نقطة الارتداد} $(\pi R, 2R)$ في
        الوعاء إلى الأسفل، قارن زمن النزول على الدويرية
        $\pi\sqrt{R/g}$ بالزمن على امتداد المزلق المستقيم
        الواصل بين النقطتين نفسيهما \emph{(بتسارع ثابت
        $g\sin\alpha$ على امتداد الوتر)}: بيّن أن الوتر يأخذ
        $\sqrt{\pi^2 + 4}\,\sqrt{R/g} \approx 3.72\sqrt{R/g}$.
        فالمنحني يغلب المستقيم --- بل هو في الواقع أسرع
        كل المنحنيات، وهي نتيجة من حساب التغيرات.
  \item بيّن أنه على القوس ($0 < t < 2\pi$)،
        \[
        \frac{\dd y}{\dd x} = \cot\frac t2,
        \qquad
        \frac{\dd^2 y}{\dd x^2} =
        -\frac{1}{4R\sin^4\frac t2} < 0 :
        \]
        فالقوس مقعّر، بمماسات عمودية عند نقاط
        الارتداد بالضبط.
  \item استعد المماس العمودي عند \omterm{ex:b1:curves:astroid}{نقطة الارتداد} في
        \cref{exo:b1:curves:2} هندسيًا: احسب الاتجاه الحدّي
        للوتر $MT$ في السؤال 3 عندما $t \to
        0^{+}$، دون أيّ نشور أصلًا.
  \item توليفة، في أربع جمل: أيّ ثلاث مبرهنات مسمّاة
        برهنت عليها هذه المسألة (بأعدادها $8R$ و $3\pi
        R^2$ و $\pi\sqrt{R/g}$ وبمؤلفيها)؛ وأيّ أداة حسابية
        واحدة (تعميلات $\sin\frac t2$ و $\cos\frac t2$)
        حرّكت الأجزاء 1 و 2 و 4 كلها؛ وكيف حوّلت
        العلاقة الذاتية $y = s^2/8R$ الهندسةَ
        إلى معادلة تفاضلية خطية؛ وعلى أيّ فصل
        من هذا الكتاب اتّكأ كل جزء.
\end{enumerate}
\end{problem}
